1: % Reference entries, B (AN-11).  Law 1: os/iword/app/face_a.c block_key
2: % (the keys that act while Block is on; "Block on", "Delete Block? No (Yes)",
3: % M_FONTBLK "Attribute: 1 Size; 2 Appearance: 0"), os/iword/core/block.c
4: % (the block from the anchor to the cursor; iw_ed_block_attr), codes.c
5: % iw_ed_attr (F6: the pair of codes, the cursor between them; F6 again moves
6: % past the off code).  Shots: ../wordperfect51-decomp mint MINT12, MINT28.
7: % GATES: test-pubs-iword test-iword-codes test-iword-face test-flair-iword-block test-flair-iword-format
8: \Gates{test-pubs-iword test-iword-codes test-iword-face test-flair-iword-block test-flair-iword-format}
9: \begin{Feature}{Block}
10: Block marks a stretch of text so that you can act on all of it at once:
11: make it bold or underlined, give it a size or an appearance, move it, copy
12: it or delete it.
13: 
14: You turn Block on with the cursor at one end of the text and move the
15: cursor to the other end.  The text between the place where you turned Block
16: on and the cursor is the block.  It is shown highlighted, and \Scr{Block on}
17: is shown at the left of the status line while Block is on.
18: 
19: To block text,
20: \begin{WPSteps}
21: \item Move the cursor to the first (or the last) character of the text.
22: \item Press \FK{Block}{Alt-F4}.  \Scr{Block on} appears on the status line.
23: \item Move the cursor to the other end of the text with the cursor keys.
24: The highlight follows the cursor.
25: \item Press the key of the feature you want to use on the block: Bold
26: (\K{F6}), Underline (\K{F8}), Font (\K{Ctrl-F8}), Move (\K{Ctrl-F4}), or
27: \K{Delete}.
28: \end{WPSteps}
29: \FlairShot{block}{640}{440}{Block on: the sentence blocked from its start to the cursor}
30: 
31: \Documents{wpkey}{Alt-F4}
32: \Documents{wpkey}{F12}
33: \Documents{wpkey}{Block:Alt-F4}
34: \Documents{wpkey}{Block:F12}
35: \Documents{wpkey}{Block:F1}
36: \Documents{wpkey}{Block:Esc}
37: \Documents{wpkey}{Block:Alt-F3}
38: \Documents{wpkey}{Block:F11}
39: \K{F12} does what Alt-F4 does: it turns Block on, and off again.  To turn
40: Block off without using it, press \FK{Block}{Alt-F4} or \K{F12} again, or
41: press \FK{Cancel}{F1} or \K{Esc}.  The text stays as it was.
42: 
43: \Notes
44: \Topic{Cursor Keys While Block Is On}
45: Every cursor key works while Block is on, including the Home combinations
46: (see \InRef{Cursor Movement}): \K{Home, Right} or \K{End} stretches the
47: block to the end of the line, \K{Home, Home, Down} to the end of the
48: document.  The block may run backward from where you turned Block on as
49: well as forward.
50: 
51: \Topic{Features That Use a Block}
52: While Block is on, InitechWord acts only on these keys: the cursor keys,
53: \K{Alt-F4}, \K{F12}, \K{F1} and \K{Esc} (Block off), \K{F6} and \K{F8} (see
54: \InRef{Bold} and \InRef{Underline}), \K{Ctrl-F8} (see \InRef{Block,
55: Attributes}), \K{Ctrl-F4} (see \InRef{Move, Block}), \K{Delete} and
56: \K{Backspace} (see \InRef{Block, Delete}), and \K{Alt-F3} or \K{F11}
57: (Reveal Codes).  Search, Replace and the other function keys do nothing
58: while Block is on.
59: Typing a character does nothing while Block is on.
60: 
61: \Topic{Highlight}
62: The block is drawn in black on grey.  The cell under the cursor is drawn in
63: reverse video, as always, and is not part of the highlight: the block ends
64: just before the cursor when the cursor is past its end.
65: 
66: \Topic{Reveal Codes}
67: With Reveal Codes on, the codes area shows \Code{[Block]} where you turned
68: Block on.  \Code{[Block]} is not a code in the document; it goes away when
69: Block is turned off.
70: 
71: \Topic{Selecting a Sentence, Paragraph or Page}
72: Instead of blocking by hand, you can have InitechWord block the sentence,
73: paragraph or page the cursor is in (see \InRef{Move, Sentence/Paragraph/Page}).
74: 
75: \SeeAlso{Block, Attributes; Block, Delete; Cancel; Move, Block; Reveal Codes}
76: \end{Feature}
77: 
78: \begin{Feature}{Block, Attributes}
79: You can give a block any appearance or size: InitechWord puts the
80: attribute's on code at the start of the block and its off code at the end,
81: and turns Block off.
82: 
83: To give a block an attribute,
84: \begin{WPSteps}
85: \item Block the text (see \InRef{Block}).
86: \item Press \FK{Bold}{F6} to make the block bold
87: \Or
88: Press \FK{Underline}{F8} to underline it
89: \Or
90: Press \FK{Font}{Ctrl-F8}.  The status line shows
91: \Scr{Attribute: 1 Size; 2 Appearance: 0}.  Select \Ch{Size}{1} or
92: \Ch{Appearance}{2}, then the size or appearance by its number, as on the
93: Size and Appearance menus.
94: \end{WPSteps}
95: \Documents{wpkey}{Block:F6}
96: \Documents{wpkey}{Block:F8}
97: \Documents{wpkey}{Block:Ctrl-F8}
98: \Documents{wpmenu}{Block:Ctrl-F8,1=Size}
99: \Documents{wpmenu}{Block:Ctrl-F8,1,1=Suprscpt}
100: \Documents{wpmenu}{Block:Ctrl-F8,1,2=Subscpt}
101: \Documents{wpmenu}{Block:Ctrl-F8,1,3=Fine}
102: \Documents{wpmenu}{Block:Ctrl-F8,1,4=Small}
103: \Documents{wpmenu}{Block:Ctrl-F8,1,5=Large}
104: \Documents{wpmenu}{Block:Ctrl-F8,1,6=Vry Large}
105: \Documents{wpmenu}{Block:Ctrl-F8,1,7=Ext Large}
106: \Documents{wpmenu}{Block:Ctrl-F8,2=Appearance}
107: \Documents{wpmenu}{Block:Ctrl-F8,2,1=Bold}
108: \Documents{wpmenu}{Block:Ctrl-F8,2,2=Undln}
109: \Documents{wpmenu}{Block:Ctrl-F8,2,3=Dbl Und}
110: \Documents{wpmenu}{Block:Ctrl-F8,2,4=Italc}
111: \Documents{wpmenu}{Block:Ctrl-F8,2,5=Outln}
112: \Documents{wpmenu}{Block:Ctrl-F8,2,6=Shadw}
113: \Documents{wpmenu}{Block:Ctrl-F8,2,7=Sm Cap}
114: \Documents{wpmenu}{Block:Ctrl-F8,2,8=Redln}
115: \Documents{wpmenu}{Block:Ctrl-F8,2,9=Stkout}
116: 
117: For example, to make the first word of a sentence bold after you have
118: typed it,
119: \BeginReplay{ref-blockattr}
120: \begin{WPSteps}
121: \item Type \TypeRun{The report is late.}
122: \item Press \PressKey{Home}, \PressKey{Left} to move to the start of the line.
123: \item Press \Press{Block}{Alt-F4}, then press \PressKey{Right} three times.\Keys{Right Right}
124: \WordScreen[rows=1]{ref-blockattr-1}{The first word blocked}
125: \item Press \Press{Bold}{F6}.  The word is bold and Block is off.
126: \item Press \Press{Reveal Codes}{Alt-F3}.
127: \end{WPSteps}
128: \WordScreen[rows=13]{ref-blockattr-2}{The block made bold: the codes at its two ends}
129: \EndReplay
130: 
131: \Notes
132: \Topic{The Cursor}
133: After the attribute is given, the cursor is at the same end of the text it
134: was on: after the off code if it was at the end of the block, after the on
135: code if it was at the start.
136: 
137: \Topic{The Font Menu With Block On}
138: With Block on, the Font menu offers only Size and Appearance.  Normal is not
139: offered.
140: 
141: \SeeAlso{Appearance; Block; Bold; Size; Underline}
142: \end{Feature}
143: 
144: \begin{Feature}{Block, Delete}
145: You can delete a whole block of text at once.
146: 
147: To delete a block,
148: \begin{WPSteps}
149: \item Block the text (see \InRef{Block}).
150: \item Press \K{Delete} or \K{Backspace}.  The status line asks
151: \Scr{Delete Block? No (Yes)}.
152: \item Type \Typed{y} to delete the block.  Any other key leaves the block as
153: it is, with Block still on.
154: \end{WPSteps}
155: \Documents{wpkey}{Block:Delete}
156: \Documents{wpkey}{Block:Backspace}
157: 
158: \Notes
159: \Topic{Codes in the Block}
160: Text outside the block keeps the appearance and size it had.  If the block
161: held the start or the end of an attribute, InitechWord puts the codes needed
162: for that where the block was.
163: 
164: \Topic{Delete With Move}
165: You can also delete a block with \FK{Move}{Ctrl-F4}: select \Ch{Block}{1},
166: then \Ch{Delete}{3} (see \InRef{Move, Block}).  Move deletes without asking.
167: 
168: \SeeAlso{Block; Delete; Move, Block}
169: \end{Feature}
170: 
171: \begin{Feature}{Bold}
172: Bold makes text stand out by printing it darker.  On the editing screen
173: bold text is drawn in bold.
174: 
175: To make text bold as you type it,
176: \begin{WPSteps}
177: \item Press \FK{Bold}{F6}.  InitechWord puts the codes \Code{[BOLD]} and
178: \Code{[bold]} in the document, with the cursor between them.
179: \item Type the text.
180: \item Press \FK{Bold}{F6} again.  The cursor moves past \Code{[bold]}, so
181: that what you type next is not bold.
182: \end{WPSteps}
183: \Documents{wpkey}{F6}
184: \begin{WPNote}To make text bold after you have typed it, block it and press
185: Bold (see \InRef{Block, Attributes}).\end{WPNote}
186: 
187: \BeginReplay{ref-bold}
188: For example, type \TypeRun{Please read the }(with the space), press
189: \Press{Bold}{F6}, type \TypeRun{new}, press \Press{Bold}{F6} again and type
190: \TypeRun{ policy.}  The screen shows:
191: \WordScreen[rows=1]{ref-bold}{A bold word typed with Bold (F6)}
192: \EndReplay
193: 
194: \Notes
195: \Topic{Bold From the Font Menu}
196: Bold is also Bold (1) on the Appearance menu of the Font key: \KC{Ctrl-F8}{2,1}
197: does what \K{F6} does.
198: 
199: \Topic{Bold in Bold Text}
200: Pressing Bold with the cursor in bold text ends the bold at the cursor: the
201: \Code{[bold]} code is moved to the cursor and the cursor goes past it.
202: 
203: \Topic{Removing Bold}
204: Delete the \Code{[BOLD]} or the \Code{[bold]} code; the other one goes with
205: it (see \InRef{Codes}).
206: 
207: \SeeAlso{Appearance; Block, Attributes; Codes; Underline}
208: \end{Feature}
